\section{Improvements made while walking through the homework}
The walkthrough was also a usability test. The following changes were implemented in the calculator package, rather than left as suggestions.

\begin{longtable}{@{}>{\raggedright\arraybackslash}p{0.21\linewidth}>{\raggedright\arraybackslash}p{0.34\linewidth}>{\raggedright\arraybackslash}p{0.37\linewidth}@{}}
\toprule Friction found & Change made & Concrete benefit\\\midrule\endhead
Repeated model entry & \texttt{LOAD + SOLVE} provides input presets for all sixteen questions, using the same general calculators. & HW5 Q6 needs eleven numeric entries manually and zero through quick load; both produce variance 56.\\
Individual marginal queries & \texttt{MARGINAL TABLES} aggregates and prints both complete finite marginal distributions. & HW4 Q1's five marginal values appear in one command instead of five separate queries.\\
An unhelpful independence-failure message & The app now displays a violating pair, its joint probability, and its marginal product. & The displayed counterexample gives a concrete starting point for the exact written non-independence proof.\\
No direct continuous marginal/conditional workflow & The new \texttt{JOINT SLICE} tool keeps the original variable names, accounts for a known global multiplier, and reports marginal values and conditional results. & HW4 Q2 and the reimbursement problem can be checked without manually rewriting every conditional density as a one-variable model.\\
Repeated joint-model entry at a new conditioning value & \texttt{NEW GIVEN VALUE} reuses the shape and multiplier, then recomputes the slice with the bounds you enter. & HW5 Q7 can be checked at $y=1/2$ and another $y$ without retyping $X+Y$.\\
Dummy-variable output on a one-variable grid & A single $y$ coordinate equal to zero is recognized as a univariate model. & HW5 Q1 no longer shows irrelevant $Y$ moments and undefined-correlation pages.\\
A stored-function tokenization defect in an optional preset & The triangular density preset stores $1-\sqrt{X^2}$ using executable calculator tokens. & The optional zero-covariance question loads correctly; this equals $1-|X|$ for real $X$.\\\bottomrule
\end{longtable}

\subsection*{Where the app saves the most work}
Finite-table arithmetic is highly automatable: normalization, all marginal masses, conditional probabilities, first/second moments, covariance, and transformed moments come from one loaded model. Given-moment questions such as HW5 Q5--Q6 are also ideal: the app expands the covariance terms and preserves their signs. The exact MGF templates handle the two assigned MGF models with minimal entry.

For joint continuous problems, the app removes most numerical integration and conditional-normalization arithmetic after the support and region have been selected. It still needs you to identify whether the region is the full support, an event, or a conditional slice. The homework's symbolic marginal formulas and independence proofs are short enough to write from the displayed setups; a numerical value cannot replace those functions or logical arguments.

\subsection*{Useful next improvements}
Useful next additions would be a probability/complement shortcut for the device question and an optional support sketch. The current tools provide numerical results and analytical templates; symbolic density formulas and proofs remain in the written solution.
